\Page{Collection OptiTAB}{Physique MPSI}{page-titre}
\thispagestyle{empty}
\vspace{9mm}
\begin{center}
\includegraphics[width=63mm]{images/logo-couleur.png}\par
\vspace{19mm}
{\otSemi\fontsize{11}{15}\selectfont\color{otGreenDark}CLASSE PRÉPARATOIRE SCIENTIFIQUE}\par
\vspace{5mm}
{\otLight\fontsize{23}{29}\selectfont Cours et exercices corrigés}\par
\vspace{10mm}
{\color{otGreen}\rule{25mm}{1.2pt}}\par
\vspace{9mm}
{\fontsize{15}{21}\selectfont Formation des images\par et optique géométrique}\par
\vspace{8mm}
{\fontsize{11}{16}\selectfont 24 exercices progressifs\par 3 défis avancés et leurs corrigés\par TP, activité numérique, DM et DS barémés}\par
\vspace{13mm}
{\sffamily\small\color{otMuted}Chapitre 01\quad /\quad Édition pilote à valider}\par
\vspace{6mm}
{\fontsize{11}{15}\selectfont Collection OptiTAB}\par
{\sffamily\small\color{otMuted}\href{https://www.optitab.net}{optitab.net}}
\end{center}

\Page{Édition et références}{Mentions éditoriales}{mentions}
\thispagestyle{mpsiFront}
\vspace{9mm}
\begin{minipage}{137mm}
\textbf{Physique MPSI}\par
Formation des images et optique géométrique\par
Chapitre 01 du projet de manuel annuel\par
\vspace{3mm}
Conception des contenus et mise en page : collection OptiTAB.\par
Charte graphique commune aux livres OptiTAB de physique-chimie : typographie, couleurs, repères de méthode et navigation entre énoncés et corrigés.\par
\vspace{3mm}
Édition pilote du 25 octobre 2026, révisée pour compléter les pages liminaires. Le cours et les entraînements du premier chapitre sont disponibles ; les autres chapitres figurent dans le plan de l'ouvrage pour préparer la suite.\par
\vspace{3mm}
\textbf{Référence pédagogique}\par
Programme de physique-chimie de la voie MPSI, annexe 2 de l'arrêté du 5 janvier 2021, NOR ESRS2035779A, publié au Bulletin officiel spécial n° 1 du 11 février 2021. Ce volume est consacré aux thèmes de physique et à la formation expérimentale de ce programme.\par
Le premier chapitre traite la section 1.1 « Formation des images ». La concordance des capacités est donnée p.\,\pageref{concordance}. L'ordre et le découpage des chapitres constituent des choix pédagogiques de la collection.\par
\vspace{3mm}
\textbf{Énoncés, figures et données}\par
Les exercices et évaluations portant la mention OptiTAB sont originaux. Les schémas sont construits dans le modèle indiqué ; une représentation schématique ou sans échelle est signalée lorsque cela est utile. Les séries créées pour l'entraînement ne sont pas présentées comme des mesures réellement réalisées.\par
Les supports authentiques extérieurs sont identifiés par leur diffuseur, leur date, la partie travaillée et leur lien d'accès. Les références et conditions de diffusion restent celles des auteurs ou organismes concernés.\par
\vspace{3mm}
© OptiTAB, 2026.\par
\href{https://www.optitab.net}{optitab.net}
\end{minipage}

\Page{Pour ouvrir le livre}{Avant-propos}{avant-propos}
\thispagestyle{mpsiFront}
\begin{minipage}{140mm}
{\fontsize{14}{20}\selectfont\itshape\color{otNavy}
Ce livre accompagne le passage du lycée à la physique de MPSI : comprendre les modèles, raisonner avec précision et apprendre à chercher de manière autonome.\par}
\vspace{4mm}
La première année de classe préparatoire demande de relier des connaissances qui pouvaient sembler séparées : une observation, un schéma, une équation, une mesure et une simulation décrivent souvent le même phénomène sous des formes différentes. L'objectif est d'apprendre à passer de l'une à l'autre en gardant le sens physique.\par
\vspace{2mm}
Chaque chapitre suit une progression commune. Un diagnostic réactive les acquis ; le cours développe les lois et leurs conditions d'application ; les méthodes montrent comment conduire une résolution. Les schémas soutiennent le raisonnement et permettent de vérifier les signes, les directions et les ordres de grandeur.\par
\vspace{2mm}
Les exercices vont des applications directes aux problèmes demandant davantage d'initiative. Les défis proposent des situations plus ouvertes, accompagnées d'indices progressifs. Les travaux pratiques et les activités numériques entraînent à confronter un modèle à des résultats et à discuter la précision d'une conclusion.\par
\vspace{2mm}
Les corrigés expliquent les choix de méthode, les calculs et l'interprétation des résultats. Ils servent à comprendre une difficulté après une recherche personnelle. Refaire un exercice quelques jours plus tard, sans regarder sa solution, reste une étape essentielle du travail.\par
\vspace{2mm}
Les DM et les DS permettent de travailler la rédaction et la gestion du temps. Des parcours sur de vrais devoirs et sujets de concours complètent l'entraînement ; les références permettent de retrouver le support d'origine et de distinguer ces documents des créations OptiTAB.\par
\vspace{2mm}
Cette édition pilote présente le premier chapitre complet et le plan du futur ouvrage. Elle permet de valider ensemble le niveau d'explication, la place des figures, la progression des exercices et la présentation avant de poursuivre. Le tableau de bord t'aide dès maintenant à repérer ce qui est acquis et ce qu'il faut reprendre.\par
\vspace{5mm}
{\sffamily\small\color{otNavy}Bon travail et bonne progression en physique.}\par
\vspace{2mm}
{\sffamily\small\color{otMuted}La collection OptiTAB}
\end{minipage}
